%-------------------------------------------------
% Wasaff Consulting — WC-2023-001
% Informe Ejecutivo — Sistema de Recuperación de Calor
% Papeles Cordillera S.A. / Planta Puente Alto
% Máximo 2 páginas. Sin ecuaciones. Orientado a decisión de inversión.
%-------------------------------------------------

\documentclass[11pt,a4paper]{article}

\usepackage[utf8]{inputenc}
\usepackage[spanish,es-tabla]{babel}
\usepackage{booktabs}
\usepackage{tabularx}
\usepackage{hyperref}
\usepackage[left=2.5cm, right=2.5cm, top=2.5cm, bottom=2.5cm]{geometry}
\usepackage{xcolor}
\usepackage{fancyhdr}
\usepackage{siunitx}
\usepackage{array}
\usepackage{parskip}

\definecolor{navyblue}{RGB}{5,28,44}
\definecolor{electricblue}{RGB}{34,81,255}
\definecolor{lightgray}{RGB}{244,244,240}

\pagestyle{fancy}
\fancyhf{}
\fancyhead[L]{\small\color{navyblue}\textbf{Wasaff Consulting} — Resumen Ejecutivo}
\fancyhead[R]{\small WC-2023-001}
\fancyfoot[C]{\small\color{gray}\thepage}
\renewcommand{\headrulewidth}{0.4pt}

\setlength{\parskip}{6pt}

%-------------------------------------------------
\begin{document}

%--- Encabezado ---
\begin{center}
  {\Large\bfseries\color{navyblue}
    Sistema de Recuperación de Calor Residual\\[0.3em]
    en Compresores de Aire Comprimido}\\[0.5em]
  {\large Planta Puente Alto · Papeles Cordillera S.A. · Grupo CMPC}\\[0.3em]
  {\small Elaborado por Wasaff Consulting para Leycero SpA · Diciembre 2023}
\end{center}

\noindent\rule{\textwidth}{1pt}

%--- Resumen en 5 líneas ---
\textbf{\color{navyblue}Situación:} Los seis compresores Kaeser de la sala de máquinas
disipan al ambiente hasta \textbf{1.212\,kW} de calor residual — energía que hoy no
genera valor. En operación normal (tres compresores activos), la pérdida es de
\textbf{632\,kW continuos}.

\textbf{\color{navyblue}Solución:} Un circuito hidráulico aislado, con bomba de circulación y
acumuladores, captura ese calor y entrega agua a \textbf{26\,°C} disponible para procesos
que hoy consumen energía primaria (vapor, gas, electricidad).

\textbf{\color{navyblue}Inversión adicional requerida:} Los módulos de recuperación Kaeser
ya están instalados en cada compresor. La obra nueva consiste únicamente en tuberías DN65/DN100,
una bomba de \textbf{0{,}55\,kW} y la instrumentación de control.

\noindent\rule{\textwidth}{0.4pt}

%--- Tabla de resultados clave ---
\begin{center}
\textbf{\color{navyblue}Resultados técnicos clave}
\end{center}

\begin{tabular}{@{}>{\bfseries}p{7.5cm}r@{}}
\toprule
Parámetro & Valor \\
\midrule
Calor recuperado — operación normal (C1+C2+C4) & 505\,kW \\
Calor recuperado — operación máxima (6 compresores) & 970\,kW \\
Temperatura de suministro agua caliente & 26\,°C \\
Temperatura entrada agua industrial fría & 6\,°C \\
Caudal circuito aislado — normal / máximo & 27\,m³/h · 52\,m³/h \\
Pérdida de carga total de la red hidráulica & 7{,}6\,kPa (normal) \\
Bomba de circulación requerida & 52\,m³/h · 2{,}4\,mca · 0{,}55\,kW \\
Volumen de acumulación instalado & 15.000\,L (10.000 + 5.000) \\
Autonomía térmica frente a paro de compresores & 57\,min \\
Tuberías de distribución & DN65 (ramales) + DN100 (colector) \\
\bottomrule
\end{tabular}

\vspace{0.4cm}

%--- Estimación de ahorro ---
\begin{center}
\textbf{\color{navyblue}Estimación de ahorro energético anual}
\end{center}

El calor recuperado reemplaza energía primaria que hoy el proceso consume. Estimando un
factor de carga de \textbf{60\,\%} anual y operación normal de 505\,kW:

\begin{tabular}{@{}>{\bfseries}p{8cm}r@{}}
\toprule
Concepto & Valor \\
\midrule
Energía térmica recuperada por año & 2.660\,MWh/año \\
Equivalente en gas natural (PCS 10\,kWh/m³) & 266.000\,m³\,GN/año \\
Ahorro estimado a \$350/m³\,GN & \$93.000.000\,CLP/año \\
\textit{(o bien, como ahorro en electricidad de calefacción eléctrica)} & \\
Equivalente eléctrico a \$120/kWh & \$319.000.000\,CLP/año \\
\bottomrule
\end{tabular}

{\small\color{gray}
  \textit{Nota: los precios de energía deben verificarse con los valores contractuales
  vigentes de Papeles Cordillera S.A. Los ahorros reales dependen del proceso destino del calor.}
}

\vspace{0.4cm}

%--- Próximos pasos ---
\begin{center}
\textbf{\color{navyblue}Próximos pasos recomendados}
\end{center}

\begin{enumerate}
  \item \textbf{Verificar $T_{\text{hot,in}}$:} medir con termopar la temperatura del agua
    en la salida del módulo de recuperación de cada compresor. Este dato ajustará el análisis
    térmico y confirmará el $UA$ requerido a los intercambiadores.
  \item \textbf{Identificar proceso destino:} cuantificar cuántos kW de calor a 26\,°C
    puede absorber el proceso y en qué horario. Esto define el \% de aprovechamiento real.
  \item \textbf{Ingeniería de detalle:} levantar isométrico de tuberías, definir accesorios
    (válvulas, filtros, expansión), cotizar bomba y automatización.
  \item \textbf{Evaluar ampliación del acumulador:} si el proceso requiere autonomía superior
    a 60\,min, aumentar de 15.000\,L a 20.000\,L (+\,5.000\,L).
\end{enumerate}

\noindent\rule{\textwidth}{1pt}
{\small\color{gray}
  Documento elaborado por Wasaff Consulting (WC-2023-001). Memoria de cálculo técnica
  completa disponible en documento adjunto. Código Python de cálculo entregado en
  repositorio del proyecto.
}

\end{document}
